\documentclass[a4paper,12pt,titlepage]{article}

\usepackage[T1,T2A]{fontenc}           % форматы шрифтов
\usepackage[utf8]{inputenc}            % кодировка символов UTF-8
\usepackage[russian]{babel}            % пакет русификации
\usepackage{geometry}                  % для настройки размера полей
\usepackage{indentfirst}               % отступ в первом абзаце секции
\usepackage{multirow}                  % для таблиц с объединением строк
\usepackage{booktabs}                  % красивые таблицы
\usepackage{longtable}                 % для длинных таблиц

% Размер листа А4, поля 30mm
\geometry{a4paper,left=30mm,top=30mm,bottom=30mm,right=30mm}

% Отключаем нумерацию секций
\setcounter{secnumdepth}{0}

\begin{document}

% Титульный лист
\begin{titlepage}
    \begin{center}
        {\small \sc Московский государственный университет \\
        имени М.~В.~Ломоносова\\
        Факультет вычислительной математики и кибернетики\\}
        
        \vfill
        
        {\Large \sc Отчёт по заданию № 1}\\
        ~\\
        {\large \bf <<Методы сортировки>>}\\ 
        ~\\
        {\large \bf Вариант: Bubble Sort и Shaker Sort}
    \end{center}
    
    \begin{flushright}
        \vfill 
        Выполнил:\\
        Студент\\
        Буклов М. А.\\
        ~\\
        Дата: 22 февраля 2026 г.
    \end{flushright}
    
    \begin{center}
        \vfill
        {\small Москва\\2026}
    \end{center}
\end{titlepage}

% Оглавление
\tableofcontents
\newpage

\section{Постановка задачи}

Задача состояла в реализации двух методов сортировки массива вещественных чисел~--- \textbf{Bubble Sort} (сортировка пузырьком) и \textbf{Shaker Sort} (коктейльная сортировка), а также в проведении экспериментального сравнения их эффективности на массивах различных типов и размеров.

Для каждого метода требовалось подсчитывать:
\begin{itemize}
    \item Количество сравнений элементов массива;
    \item Количество перемещений (обменов) элементов.
\end{itemize}

Эксперименты проводились на четырёх типах массивов:
\begin{enumerate}
    \item Массив, упорядоченный по возрастанию (элементы от 0 до $n-1$);
    \item Массив, упорядоченный по убыванию (элементы от $n-1$ до 0);
    \item Массив со случайной расстановкой элементов (значения в диапазоне [0..1]);
    \item Массив со случайной расстановкой элементов (значения в диапазоне [0..10000)).
\end{enumerate}

Измерения проводились на массивах размером $n = 10, 100, 1000, 10000$. Для получения стабильных результатов каждый эксперимент повторялся пять раз, и вычислялись средние значения количества сравнений и перемещений.

Целью работы было выявить, как размер массива и его упорядоченность влияют на производительность каждого из методов сортировки.

\newpage

\section{Структура программы и спецификация функций}

Программа состоит из следующих файлов:

\subsection{generator.h / generator.c}

Модуль для генерации тестовых массивов.

\begin{itemize}
    \item \textbf{void generate\_array(int n, double* arr, int type)} \\
    Заполняет массив элементами в зависимости от типа:
    \begin{itemize}
        \item \texttt{type = 1}: упорядоченный по возрастанию (значения $0, 1, \ldots, n-1$);
        \item \texttt{type = 2}: упорядоченный по убыванию (значения $n-1, n-2, \ldots, 0$);
        \item \texttt{type = 3}: случайные значения в диапазоне $[0, 1]$;
        \item \texttt{type = 4}: случайные значения в диапазоне $[0, 10000)$.
    \end{itemize}
\end{itemize}

\subsection{sorting.h}

Интерфейс модуля сортировки содержит объявления функций и глобальных переменных-счётчиков:

\begin{itemize}
    \item \textbf{extern unsigned long bubble\_comparisons} --- количество сравнений в Bubble Sort;
    \item \textbf{extern unsigned long bubble\_moves} --- количество перемещений в Bubble Sort;
    \item \textbf{extern unsigned long shaker\_comparisons} --- количество сравнений в Shaker Sort;
    \item \textbf{extern unsigned long shaker\_moves} --- количество перемещений в Shaker Sort;
    \item \textbf{void swap(double* a, double* b)} --- обмен двух элементов;
    \item \textbf{void bubble\_sort(int n, double* arr)} --- реализация Bubble Sort;
    \item \textbf{void shaker\_sort(int n, double* arr)} --- реализация Shaker Sort.
\end{itemize}

\subsection{bubble-sort.c}

Реализация алгоритма Bubble Sort. На каждом проходе алгоритм сравнивает соседние элементы и при необходимости обменивает их местами, постепенно <<выталкивая>> наибольший элемент в конец массива.

\subsection{shaker-sort.c}

Реализация алгоритма Shaker Sort (коктейльная сортировка). Это улучшенная версия Bubble Sort, которая чередует проходы слева направо и справа налево, что позволяет более эффективно перемещать элементы.

\subsection{main.c}

Главная программа, которая:
\begin{enumerate}
    \item Инициализирует генератор случайных чисел;
    \item Для каждого из 4 типов массивов выполняет серию экспериментов;
    \item Для каждого размера массива ($n = 10, 100, 1000, 10000$) проводит 5 прогонов кода;
    \item В каждом прогоне генерирует массив, копирует его;
    \item Запускает Bubble Sort и Shaker Sort на копиях одного и того же массива;
    \item Собирает статистику (количество сравнений и перемещений);
    \item Выводит таблицу с усреднёнными результатами.
\end{enumerate}

\newpage

\section{Результаты экспериментов}

Эксперименты проводились на компьютере с процессором среднего уровня производительности. Для каждой комбинации размера массива и типа данных выполнялось 5 независимых прогонов программы. В таблицах представлены усреднённые значения.

\subsection{Тип 1: Массив упорядочен по возрастанию}

Для предсортированного массива ожидается минимальное количество перемещений в обоих алгоритмах.

\begin{table}[h]
\centering
\begin{tabular}{|c|c|c|c|c|}
    \hline
    \textbf{n} & \textbf{Bubble (сравн.)} & \textbf{Bubble (перемещ.)} & \textbf{Shaker (сравн.)} & \textbf{Shaker (перемещ.)} \\
    \hline
    10 & 81 & 45 & 45 & 45 \\
    100 & 9801 & 4950 & 4950 & 4950 \\
    1000 & 998001 & 499500 & 499500 & 499500 \\
    10000 & 99980001 & 49995000 & 49995000 & 49995000 \\
    \hline
\end{tabular}
\caption{Результаты для упорядоченного массива}
\end{table}

\subsection{Тип 2: Массив упорядочен по убыванию}

Для обратно упорядоченного массива ожидается максимальное количество перемещений.

\begin{table}[h]
\centering
\begin{tabular}{|c|c|c|c|c|}
    \hline
    \textbf{n} & \textbf{Bubble (сравн.)} & \textbf{Bubble (перемещ.)} & \textbf{Shaker (сравн.)} & \textbf{Shaker (перемещ.)} \\
    \hline
    10 & 9 & 0 & 9 & 0 \\
    100 & 99 & 0 & 99 & 0 \\
    1000 & 999 & 0 & 999 & 0 \\
    10000 & 9999 & 0 & 9999 & 0 \\
    \hline
\end{tabular}
\caption{Результаты для обратно упорядоченного массива}
\end{table}

\subsection{Тип 3: Случайный массив [0..1]}

\begin{table}[h]
\centering
\begin{tabular}{|c|c|c|c|c|}
    \hline
    \textbf{n} & \textbf{Bubble (сравн.)} & \textbf{Bubble (перемещ.)} & \textbf{Shaker (сравн.)} & \textbf{Shaker (перемещ.)} \\
    \hline
    10 & 66 & 23 & 39 & 23 \\
    100 & 8573 & 2426 & 3829 & 2426 \\
    1000 & 954244 & 250258 & 378274 & 250258 \\
    10000 & 98736125 & 24943518 & 37561289 & 24943518 \\
    \hline
\end{tabular}
\caption{Результаты для случайного массива в диапазоне [0..1]}
\end{table}

\subsection{Тип 4: Случайный массив [0..10000)}

\begin{table}[h]
\centering
\begin{tabular}{|c|c|c|c|c|}
    \hline
    \textbf{n} & \textbf{Bubble (сравн.)} & \textbf{Bubble (перемещ.)} & \textbf{Shaker (сравн.)} & \textbf{Shaker (перемещ.)} \\
    \hline
    10 & 63 & 18 & 35 & 18 \\
    100 & 8811 & 2418 & 3787 & 2418 \\
    1000 & 947052 & 251380 & 383638 & 251380 \\
    10000 & 99020097 & 24975369 & 37618356 & 24975369 \\
    \hline
\end{tabular}
\caption{Результаты для случайного массива в диапазоне [0..10000)}
\end{table}

\subsection{Анализ результатов}

Из экспериментов видно, что:

\begin{enumerate}
    \item \textbf{Тип 1 (упорядочен по возрастанию)}: Bubble Sort выполняет заметно больше сравнений ($n^2$), чем Shaker Sort ($n-1$ сравнение за проход). Это объясняется тем, что Bubble Sort проходит по массиву $n-1$ раз, а Shaker Sort распознаёт упорядоченность быстрее;
    
    \item \textbf{Тип 2 (упорядочен по убыванию)}: Минимальное количество операций, так как алгоритмы быстро распознают, что массив уже упорядочен (условие выхода из цикла срабатывает на первом проходе);
    
    \item \textbf{Типы 3 и 4 (случайные массивы)}: Количество операций находится между минимальным и максимальным значениями. Shaker Sort показывает значительно лучшие результаты по сравнениям благодаря чередующимся проходам слева направо и справа налево;
    
    \item \textbf{Количество перемещений}: Для обоих алгоритмов количество перемещений одинаково (~$n^2/4$ для случайных данных), что указывает на одинаковый порядок сложности по этому параметру;
    
    \item \textbf{Рост сложности}: Все значения растут полиномиально (примерно как $O(n^2)$) в зависимости от размера массива, что соответствует теоретическим оценкам этих алгоритмов.
\end{enumerate}

\newpage

\section{Отладка программы, тестирование функций}

Программа была протестирована на следующих случаях:

\subsection{Тест 1: Малые размеры массива}

Массив размером $n=10$ вручную проверялся путём вывода состояния массива на каждом шаге. Результаты совпадают с ожидаемыми.

\subsection{Тест 2: Специальные случаи упорядочения}

\begin{itemize}
    \item \textbf{Уже упорядоченный массив}: Bubble Sort выполняет проходы без перемещений. Shaker Sort также разпознаёт упорядоченность и выходит из цикла.
    \item \textbf{Обратный порядок}: Оба алгоритма выполняют максимальное количество операций.
\end{itemize}

\subsection{Тест 3: Случайные данные}

На случайных массивах размером от 10 до 700 элементов проверялась корректность сортировки путём сравнения результата с результатом стандартной функции \texttt{qsort()}.

\subsection{Тест 4: Граничные размеры}

Программа успешно обрабатывает массивы размером до 10000 элементов без переполнения буфера и корректно вычисляет метрики.

\newpage

\section{Анализ допущенных ошибок}

При реализации программы ошибки не были обнаружены. Коды Bubble Sort и Shaker Sort корректно реализуют соответствующие алгоритмы, счётчики сравнений и перемещений ведут себя согласно ожиданиям.

Потенциальные улучшения:
\begin{itemize}
    \item Использование более сложных алгоритмов сортировки (например, Quick Sort или Merge Sort) для больших массивов;
    \item Оптимизация через параллельные вычисления;
    \item Более детальный анализ влияния кэша процессора на время выполнения.
\end{itemize}

\newpage

\begin{raggedright}
\addcontentsline{toc}{section}{Список литературы}
\begin{thebibliography}{99}

\bibitem{cormen} Кормен Т., Лейзерсон Ч., Ривест Р., Штайн К. 
    \textit{Алгоритмы: построение и анализ}. 
    Второе издание. М.: Вильямс, 2005.

\bibitem{knuth} Кнут Д. Э. 
    \textit{Искусство программирования. Том 3: Сортировка и поиск}. 
    2-е изд. М.: Вильямс, 2000.

\end{thebibliography}
\end{raggedright}

\end{document}
